\scanpage{361}
\mspage{100}

se factorise (évidemment de façon unique) en un foncteur
\[
H_C:T_0^\circ\longrightarrow\mathrm{Cat\,hot}
\]
i.e. qu'on a \deleted{\unsure{donc}} un diagramme commutatif
\diagram{\begin{tikzcd}[ampersand replacement=\&]
T^\circ\arrow[r,"{t\mapsto C_t}"]\arrow[d]\&\mathrm{Cat}\arrow[d]\\
T_0^\circ\arrow[r,"{\substack{H_C\\\deleted{\unsure{C}}}}"]\&\mathrm{Cat\,hot}.
\end{tikzcd}}

Nous allons supposer de plus

\margin{(Hypothèse de filtration)}L'ens ordonné $T_0$ associé à $T$ est \uline{filtrant} décroissant \inserted{(filtrant à gauche)}, i.e. si $t,t'\in\mathrm{Ob}\,T$, il existe $t''\in\mathrm{Ob}\,T$ \deleted{\unsure{tel que}} des flèches $t^{\unsure{\prime\prime|\prime\prime\prime}}\mathrel{\unsure{-|\to}}t$, $t''\mathrel{-}t'$.

Cette hypothèse est bien satisfaite dans $\underline{\mathrm{Ty}}$, car pour $\tau,\tau'\in\mathrm{Ob}\,\underline{\mathrm{Ty}}=\mathrm{Ty}$, il suffit de prendre $\tau\tau'$, qui s'envoie dans $\tau$ par $\mathrm{id}_\tau*p_{\tau'}$, dans $\tau'$ par $p_\tau*\mathrm{id}_{\tau'}$, où $p_\tau:\tau\longrightarrow1$ est le morphisme de $\tau$ dans l'objet final. \deleted{\unsure{L|E}}

Ainsi, $\deleted{\unsure{\Phi|H}}H_C$ définit un Ind-objet de $\mathrm{Cat\,hot}$, \inserted{que nous allons désigner par $\underline{\underline{\unsure{\mathcal C|C}}}$}. Si on suppose de plus

\margin{Hypothèse de dénombrabilité (tjrs satisfaite si $\mathrm{card}\,T_0\leq\aleph_0$)}$T_0$ admet un sous-ensemble \uline{dénombrable} cofinal à gauche,

alors cet ind-objet de $\mathrm{Cat\,hot}$ peut être défini, \unsure{\deleted{de}|de} \unsure{beaucoup|\deleted{beaucoup}|\deleted{la manière}} de \unsure{façons|façon}\unsure{,} en termes

% archive sheet number: 360
